\documentclass[11pt,a4paper]{article}

\usepackage[utf8]{inputenc}
\usepackage[T1]{fontenc}
\usepackage{geometry}
\usepackage{hyperref}
\usepackage{booktabs}

\geometry{margin=2.5cm}
\hypersetup{colorlinks=true, linkcolor=blue, urlcolor=blue, citecolor=blue}

\title{Symbolic Provenance-Gated Defense against Indirect Prompt Injection\\[0.5em]\large Pre-Final Project Progress Report}
\author{Murtazin Aizat}
\date{July 17, 2026}

\begin{document}
\maketitle
\sloppy

\noindent\textbf{Code repository:}\\
\url{https://github.com/muitiiifruckt/neural-taint-tracking-thesis}\\[0.4em]
\textbf{Project site (diagrams, full results):}\\
\url{https://muitiiifruckt.github.io/neural-taint-tracking-thesis/}

\section*{Summary}

Following reviewer feedback, the project is narrowed and reframed as a \textbf{symbolic
provenance-gated defense} (not ``neural taint tracking''), focused on stages T1--T3. The
research question is: \emph{does field-sensitive provenance gating (T3) reduce targeted attack
success rate (ASR) while preserving benign utility better than a strict gate (T2)?} A learned
sanitizer and a formal noninterference proof are explicitly out of scope.

Since the previous stage, all four treatments T0--T3 are now fully implemented, unit-tested,
and evaluated end-to-end on the AgentDojo benchmark. This report presents the complete T0--T3
comparison and answers the research question.

\section{What Has Been Done So Far}

\paragraph{Repository and reproducibility.} All code is committed and pushed (the previous
submission's gap): \texttt{src/taint/} (taint core + gate policy + AgentDojo pipeline
integration), \texttt{src/agent/} (pipeline wrapper + shared metric arithmetic),
\texttt{tests/} (29 unit tests, no network needed), \texttt{scripts/}, \texttt{docs/}
(per-stage specifications), and \texttt{results/} (summaries + per-case taint traces).
Dependencies are pinned; the evaluation runner asserts the expected case counts on every full
run. A themed documentation site presents the diagrams and results.

\paragraph{Three separated evaluations.} The benchmark reporting is split into three
evaluations that must never be mixed into one number:
\begin{enumerate}
  \item \emph{Benign user-task evaluation} --- 10 user tasks, no attack.
  \item \emph{Direct injection-goal achievability} --- the 14 injection goals executed
        directly as ordinary requests: \textbf{7 of 14} are achievable when simply asked.
  \item \emph{Indirect prompt-injection security} --- 140 attacked cases. The
        \textbf{primary} ASR metric is computed on the direct-achievable subset (70 cases);
        the full-set ASR (140 cases) is secondary, since the full set mixes in 7 goals the
        agent cannot perform anyway (attacks on them succeed in 0/70 cases).
\end{enumerate}

\paragraph{Corrected metric semantics.} While wiring the T2 gate, a sign error was found in
the metric aggregation: AgentDojo's raw per-case \texttt{security} field is \texttt{True} when
the \emph{injection goal was achieved} (confirmed by the \texttt{BaseInjectionTask.security}
docstring, by \texttt{TaskSuite.check} which runs the attack's ground truth and requires
\texttt{True}, and empirically on the logs). Earlier scripts read it inverted, so all
previously reported ASR values were the complement of the true ones (e.g.\ the earlier
``75.7\% ASR'' is in fact a 75.7\% \emph{security rate}, i.e.\ 24.3\% ASR). The arithmetic is
now centralized in \texttt{src/agent/metrics.py} and guarded by a regression test that
reproduces exactly this inversion. All numbers below use the corrected semantics.

\paragraph{T1 --- tag-only taint tracking.} A forked tools executor
(\texttt{ObservingToolsExecutor}) is integrated into a live AgentDojo pipeline; it labels the
\emph{raw structured} tool output before it is flattened to a string, preserving per-field
provenance. Real traces (150 files) show untrusted tool outputs propagating into later
tool-call arguments with parent lineage. Running the real benchmark also surfaced and fixed a
genuine matching bug (short values extracted from a larger untrusted block were not detected).

\paragraph{T2 / T3 --- the gate.} \texttt{src/taint/policy.py} defines an explicit list of 10
privileged (state-changing) tools and a per-tool table of sensitive fields (routing /
exfiltration targets and destructive-target ids are sensitive; content fields are not).
\texttt{GatingToolsExecutor} blocks a privileged call --- the call is \emph{not executed}, the
model receives an explanatory tool error, and the trace records the block reason and the exact
tainted argument paths. T2 blocks on any tainted argument; T3 blocks only on tainted sensitive
fields.

\section{Results: T0--T3 Comparison}

All numbers are on the same 150 cases (10 benign + 140 security), same model
(\texttt{gpt-4o-mini-2024-07-18}), suite \texttt{workspace} (v1.2.2), attack
\texttt{important\_instructions}.

\begin{table}[h]
  \centering
  \caption{Metric comparison across stages. Targeted ASR (subset) is the primary metric.}
  \begin{tabular}{@{}lcccc@{}}
    \toprule
    \textbf{Metric} & \textbf{T0} & \textbf{T1} & \textbf{T2 (strict)} & \textbf{T3 (field)} \\
    \midrule
    Benign Utility & 100.0\% & 100.0\% & 80.0\% & 80.0\% \\
    Utility Under Attack & 20.0\% & 20.0\% & 34.3\% & 36.4\% \\
    \textbf{Targeted ASR, subset (primary)} & \textbf{48.6\%} & \textbf{51.4\%} & \textbf{11.4\%} & \textbf{8.6\%} \\
    Targeted ASR, full set & 24.3\% & 28.6\% & 5.7\% & 4.3\% \\
    Benign overblock rate & --- & 0\% & 20.0\% & 20.0\% \\
    Blocked calls (total) & --- & 0 & 579 & 496 \\
    Taint leakage rate & --- & 36.8\% & 0.0\% & 2.1\% \\
    \bottomrule
  \end{tabular}
\end{table}

\paragraph{T0 $\to$ T1.} Benign Utility and Utility Under Attack are \emph{identical} to T0 ---
the acceptance criterion for a purely observational layer. The small ASR difference is
run-to-run LLM noise; the gate is disabled and cannot affect outcomes. Taint leakage of 36.8\%
means over a third of executed privileged calls carried untrusted content --- the direct
motivation for a gate.

\paragraph{T1 $\to$ T2.} The strict gate cuts subset ASR almost $4.5\times$ (48.6\% to
11.4\%). The cost is Benign Utility dropping from 100\% to 80\%: two benign tasks break because
they legitimately copy a meeting participant's email address from an earlier message into
\texttt{create\_calendar\_event.participants} --- the same data pattern the attack uses, but
with no malicious intent.

\paragraph{T2 $\to$ T3 (the research question).} Field-sensitive gating lowers subset ASR
further (11.4\% to 8.6\%) and raises Utility Under Attack (34.3\% to 36.4\%), because it stops
blocking calls whose only tainted argument is a \emph{content} field (e.g.\ an email body
quoting injected text) --- 496 blocked calls vs.\ 579 under T2. \textbf{However, Benign Utility
is not recovered:} both policies break the same two benign tasks, because the failures live in
the \texttt{participants} field, which is sensitive under both policies (it is the same field
an attacker uses to redirect invites). So the trade-off between T2 and T3 is close to
degenerate on this suite: T3 is marginally safer and marginally more useful under attack, but
no less intrusive to the benign user. The bottleneck is precision on \emph{routing} fields, not
content fields --- a limitation a purely provenance-based gate cannot resolve.

\paragraph{Residual attacks.} Every attack that survives T2 (8/140) and T3 (6/140) is the same
goal: deleting a file by a short numeric id. The cause is in T1, not the gate --- the id is
shorter than the 4-character minimum-match threshold, so the taint layer never labels it
untrusted and both gates see a trusted argument. This is a documented, conscious trade-off (the
threshold prevents false matches on short strings), not a hidden gap.

\section{Progress Compared to the Plan}

\begin{table}[h]
  \centering
  \begin{tabular}{@{}lll@{}}
    \toprule
    \textbf{Stage} & \textbf{Description} & \textbf{Status} \\
    \midrule
    T0 & Baseline without defense & Completed \\
    T1 & Tag-only taint tracking & Completed (full eval, traces, leakage metric) \\
    T2 & Strict gate & Completed (full eval, false-positive analysis) \\
    T3 & Field-sensitive gate & Completed (full eval, answers research question) \\
    T4 & Rule-based sanitizer & Optional --- not attempted \\
    \bottomrule
  \end{tabular}
\end{table}

The previous report claimed T1 was $\sim$60--70\% done, and the code was not verifiable in the
repository. Since then: the repository is complete and pushed; T1, T2, and T3 are each 100\%
done including full 150-case evaluations; T0 reporting is separated into the three evaluations;
the metric-semantics bug is fixed and regression-tested; and the missing \texttt{src/agent/}
directory was added. All ten artifacts on the reviewer's checklist are delivered: clean
repository layout, real T1 integration, real traces, T1 non-regression vs.\ T0, taint leakage
metric, T2 strict gate with a privileged-tool list, T3 field-sensitive gate with a
sensitive-fields table, the T0--T3 comparison with primary/secondary ASR and overblock rate,
unit tests (propagation / strict / field / metric non-regression), and reproducibility
(pinned dependencies, exact commands, model, benchmark version, task list, case-count asserts).

\section{Plan for the Final Report}

The technical ladder T0--T3 is complete; the remaining work is the final write-up and, if time
allows, one optional extension:

\begin{enumerate}
  \item \textbf{Final report text:} idea $\to$ threat model $\to$ per-stage results $\to$
        research-question answer $\to$ limitations $\to$ future work. Most of this content
        already exists in the results summaries and the project site and needs consolidation.
  \item \textbf{Limitations / future work} (documented, not hidden): (a) the short-id bypass of
        the minimum-match threshold, potentially closable with provenance-aware exact-field
        matching instead of substring matching; (b) the \texttt{participants}/routing false
        positives, which need a policy beyond provenance (e.g.\ an address-book allowlist).
  \item \textbf{Optional T4:} a lightweight rule-based sanitizer or allowlist for
        routing fields, attempted only if time remains --- it is out of the required scope.
\end{enumerate}

\section{Conclusion}

The project now has a reproducible T0--T3 evaluation on a single model, a full metric
comparison, and documented false positives and residual attacks. The strict gate (T2) sharply
reduces attack success at a measurable utility cost; the field-sensitive gate (T3) is slightly
safer but does not recover that cost, because the remaining benign failures are concentrated in
routing fields that both policies must treat as sensitive. This honest answer --- rather than
``T3 is strictly better'' --- is the main substantive result, and it sets up the final report's
discussion of what a provenance-only gate can and cannot achieve.

\end{document}
